\section{Programska implementacija}
\label{sec:implementacija}

Sustav je implementiran u programskom paketu MATLAB. Vodeće načelo implementacije
bilo je ponovno korištenje postojećih, provjerenih komponenti i oslanjanje na
ugrađene MATLAB funkcije, čime je količina novoga koda svedena na najmanju mjeru.

\subsection{Ponovno korištene komponente}

Korak 2D korelacije po kameri (Poglavlje \ref{sec:metoda}) proveden je postojećim
algoritmom globalne korelacije konačnih elemenata. Rezultat -- polja čvornih 2D
pomaka za obje kamere kroz svih 300 obrađenih deformacijskih stanja -- korišten je
kao ulaz u stereo rekonstrukciju.

\subsection{Nove komponente}

Dodane su tri kratke programske jedinice, prikazane u Tablici \ref{tab:moduli}.
Triangulacija je izvedena metodom DLT preko ugrađene funkcije za dekompoziciju na
singularne vrijednosti (\texttt{svd}), a vizualizacija ugrađenim funkcijama
\texttt{trisurf} i \texttt{patch}.

\begin{table}[H]
  \centering
  \caption{Nove programske jedinice modula \texttt{Stereo3DDIC}.}
  \label{tab:moduli}
  \begin{tabularx}{\textwidth}{lX}
    \toprule
    \textbf{Datoteka} & \textbf{Uloga} \\
    \midrule
    \texttt{stereoTriangulate.m} & Linearna (DLT) triangulacija točaka iz dviju
      kamera prema \eqref{eq:dlt}. \\
    \texttt{surfaceStrains.m} & Green--Lagrangeove površinske deformacije po elementu
      prema \eqref{eq:green-lagrange}. \\
    \texttt{StereoDIC3D.m} & Glavna skripta: učitavanje podataka, triangulacija po
      frameovima, deformacije, spremanje i vizualizacija. \\
    \bottomrule
  \end{tabularx}
\end{table}

\subsection{Provjera ispravnosti}

Ispravnost implementacije potvrđena je na dva načina. Prvo, triangulacija referentnih
(nedeformiranih) položaja rekonstruira zadanu 3D mrežu s pogreškom reda veličine
$10^{-15}$, tj. do strojne preciznosti, čime je potvrđena ispravnost triangulacijskog
postupka i orijentacije projekcijskih matrica. Drugo, srednja reprojekcijska pogreška
trianguliranih deformiranih točaka iznosi približno $0{,}2$ piksela (subpikselno), što
potvrđuje ispravno pridruživanje komponenti pomaka i visoku kvalitetu korelacije.
